\documentclass[11pt]{article}
\usepackage[T1]{fontenc}
\usepackage{lmodern,amsmath,amssymb,amsthm,amscd,mathtools,mathrsfs}
\usepackage[a4paper,margin=28mm]{geometry}
\usepackage[colorlinks=true,linkcolor=blue,urlcolor=blue]{hyperref}
\setlength{\emergencystretch}{2em}
\newtheorem{proposition}{Proposition}[section]
\newtheorem{theorem}[proposition]{Theorem}
\newtheorem{lemma}[proposition]{Lemma}
\newtheorem{corollary}[proposition]{Corollary}
\theoremstyle{definition}
\newtheorem{definition}[proposition]{Definition}
\theoremstyle{remark}
\newtheorem{remark}[proposition]{Remark}
\newtheorem{example}[proposition]{Example}
\newcommand{\Q}{\mathbf Q_\ell}
\newcommand{\RHom}{R\operatorname{Hom}}
\newcommand{\Perf}{\operatorname{Perf}}
\newcommand{\Sk}{\mathbf{Sk}}
\newtheorem{problem}[proposition]{Problem}
\expandafter\def\expandafter\UrlBreaks\expandafter{\UrlBreaks\do\-}
\newcommand{\mainref}[1]{\textnormal{[Nearby, \nolinkurl{#1}]}}
\title{Exceptional and Completed Root Kernel Calculus\\
\large A Single-Transition Problem}
\author{YINGFENG JIANG}
\date{26 September 2026 --- research draft; the kernel calculus is not constructed}
\begin{document}
\maketitle
\begin{abstract}
This research draft separates the unconstructed global root shriek
interface from the nearby kernel descent theorem.  Ordinary stalkwise
inverse gluing of the distribution diagonal is impossible.  The actual
Verdier mate gives a dual-unit trace system with the full dualizing
complex, which need not be invertible.  We formulate one decisive
unit--convolution--localization problem on one complete nodal support
and one actual transition.  The recorded local calculations and finite
diagrams do not solve that problem.  No global realization theorem is
claimed.
\end{abstract}

\section{The single-transition problem}
Retain $xy=t^4$, $B=\{2\}$ and the four fixed integer-cut models of
\cite{Nearby}.  First choose the complete support
$X=u_{\{1,2\},B}^{-1}(E_B)$, containing the newly introduced component
and its node, and one actual neutralized root-level transition over
this support.  Retain finite coefficients $\Lambda_r$, independent
continuous-action and distribution indices, then adic recovery and
only then rationalization.  No general-root statement is sought.

\begin{problem}[A unit calculus on one complete nodal support]
\label{problem:unit-calculus}
\leavevmode\par\noindent
Construct an exceptional/completed or mixed kernel category on this
support, with actual continuous mapping spectra, an actual trace
transition, a specified composition $\star$ and a unit $I_X$.
On at least one rectangle crossing from an open branch to its node
(or the reverse), construct and prove the equivalences
\[
 I_X\star P\simeq P,\qquad P\star I_X\simeq P,
\]
including their naturality and unit identities.  Retain the actual
open/closed localization triangle and its connecting morphism, and
prove that the specified transition and unit comparisons respect it.
Neither quotienting kernel maps by their actions nor defining the
category by the desired equivalences counts as a construction.
\end{problem}

A solution must specify the finite predecessors, the variance of every
transition, the order of completion, the mapping objects, and the
geometric tensor/Hom and exceptional-pullback operations used to define
composition.  The actual bar or correspondence maps must prove both
unit laws on the crossing kernel.  If $I_X$ is naturally
$\Delta_*\omega_X$, its composition must be proved unital with that
object.  This is a different presentation until an additional comparison
with ordinary shriek kernels and their unit has been established.
Renaming this dual-unit does not supply such a comparison.

The ordinary-pullback obstruction is \mainref{cor:root-diagonal-stalk-obstruction}:
the degree-zero attaching square would require $\iota c=1$, impossible
already modulo $\ell$.  The actual mate is instead
\[
 (q\times q)_!\Delta_*\omega_Y\longrightarrow\Delta_*\omega_X.
\]
The node calculation \mainref{lem:root-nodal-dualizing-correction}
gives $H^{-2}=\Lambda_r(1)^2$ and $H^{-1}=\Lambda_r$ on a nodal atlas.
Thus an invertible scalar shift cannot identify this with the ordinary
unit.  These are two distinct obstacles, not missing higher homotopies.

\section{Dependency boundary and retained evidence}
The nearby closure, one-sided descent and Morita--Verdier realization
of \cite{Nearby} remain independent proved manuscript statements on
their specified domain.  Its completed residual-gerbe unit, normalized
rectangles, two coarse convolutions and carrier comparison are local
inputs only.  Full-fibre generators and integration remain unproved.
The point-valued unit/rectangle source, ambient inertia graph objects,
and actual-open geometric passage tests in the 27 September revision
of \cite{Nearby} are three distinct domains.  In particular the graph
calculation is not membership in that generated source, and the two
$U=\mathbf G_m$ tests do not supply a common strong realization.
References to their diagrams below are restricted accordingly.
First solve Problem~\ref{problem:unit-calculus}; only then address
full-fibre realization and strong convolution, and subsequently the
four-model Verdier/refinement diamond.  The diagrams recorded below
are supporting material for that later stage, not a substitute for
the missing kernel category or its unit laws.

References of the form [Nearby, label] identify exact statements in
\nolinkurl{verdier_kernel_skeletonization.tex}; they do not declare new
assumptions or import a global root interface.  The following sections
preserve earlier proved local calculations and conditional diagram
arguments that have been removed from the main manuscript.

\section{Continuous local precursor}
Use the neutralized residual groups $\Gamma_a=0$ or $\mathbf Z_\ell(1)$,
$G_n=\Gamma/\ell^n\Gamma$ and $\Lambda_r=\mathbf Z/\ell^r$.
\begin{definition}[The continuous local kernel categories]
\label{def:continuous-local-root}
At fixed $r$, use the unbounded derived category of discrete
continuous $\Gamma_a\times\Gamma_b$-modules over $\Lambda_r$, where
$\Gamma_a,\Gamma_b$ are either $0$ or $\Gamma$.  Morphisms are its
full derived mapping spectra.  These are actual sheaves on the
relative classifying topoi; no quotient by coefficient actions is
made.  Put
\[
 A_{a,r}=C^*_{\mathrm{cts}}(\Gamma_a,\Lambda_r),\qquad
 I_{a,r}=C_{\mathrm{lc}}(\Gamma_a,\Lambda_r),
\]
where $(g,h)$ acts on $I_{a,r}$ by translation through $h-g$.
Define the cohomological composition by
\[
 L\mathbin{\circledast_c}P
  =R\Gamma_{\mathrm{cts}}\bigl(\Gamma_b,
       p_{ab}^*P\otimes_{\Lambda_r}^L p_{bc}^*L\bigr).
\]
The outer group actions remain.  Take the full stable idempotent
closure generated by the $I_{a,r}$ and the constant rectangular
kernels $\Lambda_r$ on every pair, under this composition and tensor
with perfect $\Lambda_r$-complexes.  This definition allows the
infinite-rank discrete modules $I_{a,r}$.
\end{definition}
The notation $\circledast_c$ deliberately does not denote a root
shriek operation.  Continuous group cochains give its functor on
actual morphisms and Fubini gives its enhanced associators.

\begin{theorem}[Continuous diagonal and cohomological kernel comparison]
\label{thm:continuous-root-diagonal}
The actual finite root diagonal kernels have a canonical inflation
system with colimit $I_{a,r}$, which is the unit for
$\circledast_c$.  The functor
\[
 \mathcal R_{c,r}(P)=
 R\Gamma_{\mathrm{cts}}(\Gamma_a\times\Gamma_b,P)
\]
with its two cochain algebra actions is defined on full kernel
mapping spectra.  On the generated categories above it has
canonical equivalences
\begin{align*}
 \mathcal R_{c,r}(I_{a,r})&\simeq A_{a,r},\\
 \mathcal R_{c,r}(\Lambda_r)&\simeq A_{a,r}\otimes_{\Lambda_r}A_{b,r},\\
 \mathcal R_{c,r}(P)\otimes_{A_{b,r}}^L\mathcal R_{c,r}(L)
      &\xrightarrow{\sim}\mathcal R_{c,r}(L\circledast_c P).
\end{align*}
The first is an equivalence of unit \emph{bimodules}.  These
constraints preserve units and composition on the actual mapping
categories.  On the integral cartesian systems generated by these
kernels the comparisons pass through adic recovery and then
rationalization.  At the node the first comparison thereby identifies
the continuous root diagonal with the regular nearby bimodule $H_q$.
\end{theorem}
\begin{proof}
The representable finite diagonal of $BG_n$ is finite etale.
Its actual shriek kernel is the permutation module
\[
 I_{n,r}=\operatorname{Ind}_{\Delta G_n}^{G_n^2}\Lambda_r
          =\operatorname{Map}(G_n,\Lambda_r),
 \qquad \operatorname{rank}_{\Lambda_r} I_{n,r}=\ell^n.
\]
For $m\geq n$ pullback of functions along $G_m\to G_n$ defines
an equivariant map from the inflated $I_{n,r}$ to $I_{m,r}$.
Their colimit is $I_r=C_{\mathrm{lc}}(\Gamma,\Lambda_r)$.
This is the continuous coinduced module from $\Delta\Gamma$;
the explicit complement to that diagonal makes coinduction exact
(it is a filtered colimit of finite products).  Its continuous bar
Shapiro map, obtained by restriction and evaluation at the identity
coset, gives
\[
 R\Gamma_{\mathrm{cts}}(\Gamma^2,I_r)
       \xrightarrow{\sim}R\Gamma_{\mathrm{cts}}(\Delta\Gamma,\Lambda_r).
\]
The two projection actions of $A_{q,r}$ restrict to the two identical
identity maps of $\Delta\Gamma\simeq\Gamma$.  Thus this is the
regular bimodule comparison, including both actions and the module
map itself.  It is not merely a comparison of cohomology dimensions.
Equivalently compute on finite bar nerves with $I_{n,r}$ and pass
to inflation: every continuous cochain with these discrete values
factors through one finite level.  This justifies the colimit on
the actual growing kernels independently of the selected-coefficient
theorem of \cite{Local}.

The same Shapiro map with coefficients proves
$I_b\circledast_c P=P$ and $P\circledast_c I_a=P$.
Indeed the tensor with the function module is the corresponding
coinduced coefficient, and evaluation at the identity coset is the
unit constraint.  The constant--constant composite is
$A_{b,r}\otimes\Lambda_r$ on the outer group product.  The complex
$A_{b,r}$ is perfect over $\Lambda_r$: the two-term augmentation
resolution of $\Lambda_r$ over $\Lambda_r[[\Gamma_b]]$ computes it.
These calculations give closure on the identity/rectangle family
and then on its full thick closure, with all maps retained.

The composition comparison in the theorem is cup product after
pullback to the three group factors, balanced over $A_{b,r}$,
followed by continuous-cochain Fubini.  For two rectangles it is
the Kunneth equivalence computed by the tensor product of the
finite augmentation resolutions.  With an identity factor it is
the specified Shapiro/unit map.  Both sides are exact separately
in $P,L$ and preserve retracts, so these checks prove the comparison
on their thick closures; the preceding family calculation already
shows that convolution adds no other types of generators.
Naturality and composition coherence are those of cup product and
Shapiro on the same group bar nerves, not choices on isomorphism
classes.  This also proves that the map is defined on all derived
kernel morphisms, without claiming full faithfulness.

The basic kernels are flat over $\Lambda_r$ and their reductions
are the stated basic kernels at the previous precision.  Form the
full finite stable/retract/convolution closure inside the category
of derived cartesian systems in $r$.  The cochain outputs of the
basic kernels are the finite augmentation complexes above, which
commute with derived reduction; exactness and retracts extend this
to these generated systems.  Thus the displayed equivalences pass
through $R\varprojlim_r$ in the derived adic category.  Only then
invert $\ell$.  The actual local star comparison of \cite{Local}
identifies $C^*_{\mathrm{cts}}(\Gamma,\mathbf Z_\ell)[1/\ell]$
with $H_q$ as an algebra, so the regular bimodule comparison becomes
the asserted nearby unit comparison.  On geometric supports this
is the local node coefficient of $\Delta_*H$, not a proof of its
global attachment across every model.
\end{proof}

\begin{proposition}[The dual of the continuous diagonal is completed]
\label{prop:root-diagonal-completion}
The scalar dual of $I_r$ is the pseudocompact module
\[
 I_r^\vee=\varprojlim_n\Lambda_r[G_n],
\]
with transition maps summing coefficients over quotient fibres.
Continuous scalar biduality recovers $I_r$; ordinary discrete
coefficient duality does not keep it in the same finite-coefficient
class used in \cite{Local}.
\end{proposition}
\begin{proof}
The dual of the pullback map of finite function modules is the
quotient-fibre sum on distributions.  Derived Hom takes their
filtered colimit to the inverse limit of their duals.  These inverse
maps are surjective, so there is no higher inverse-limit term.
A continuous functional on the resulting profinite module with
values in the finite discrete ring factors through a finite quotient.
Its dual is therefore the original colimit of functions.  The
pairing is evaluation, equivariant for the two contragredient
actions.  Recovering $r$ gives the completed integral distribution
module; rationalization must still follow this recovery.  This
constructs the required local dual pair, but does not supply a
global six-operation theory mixing all discrete and completed kernels.
\end{proof}


\section{Coefficient branch and later coherence material}
\subsection{An actual branch calculation and its dual boundary}
\label{sec:root-branch-calculation}
The following calculation supplies a coefficient branch face of
the missing four-model gluing diagram.  It deliberately uses the
normal-pair Kummer map, not central-gerbe purity.

\begin{lemma}[Continuous Kummer link and its trace-dual boundary]
\label{lem:root-branch-link}
Fix one branch coordinate $x$ in one of the four models.  On its
strictly local smooth branch let $j_n$ be the punctured inclusion
and $i_n$ the closed inclusion in the Kummer chart
$x=z_n^{\ell^n}$.  Set
$L_{n,r}=i_n^*Rj_{n,*}\Lambda_r$, with its actual deck action.
The pullback comparisons for $z_n=z_m^{\ell^{m-n}}$ give
\[
 \operatorname*{hocolim}_n L_{n,r}\simeq E_r
\]
as a discrete continuous inertia object, through the actual unit.
For a selected unipotent branch coefficient $V$, the corresponding
continuous link is its continuous inertia module, with the sign
of the specified branch coordinate.  These are natural comparisons
on the selected mapping category.  Dualizing the finite localization
diagrams gives their actual trace-dual inverse diagrams and
connecting arrows.
\end{lemma}
\begin{proof}
The punctured strict branch has cohomology $\Lambda_r$ in degree
zero and $\Lambda_r(-1)$ in degree one, and no higher cohomology.
Kummer functoriality sends the degree-one class of $z_n$ to
$\ell^{m-n}$ times that of $z_m$; it acts identically on degree zero.
These statements identify the maps of the actual link complexes,
not an arbitrarily selected splitting.  The maps are induced by
the same Kummer morphisms as the deck actions.  At fixed $r$ the
positive-degree system is zero after sufficiently many transitions.
Filtered colimits in discrete continuous modules are exact, so the
actual unit $E_r\to\operatorname*{hocolim}_n L_{n,r}$ is an
equivalence; its degree-zero action is the trivial action.
For a selected coefficient use its finite-action presentation and
pull back to a Kummer level trivializing it.  The same calculation
with that coefficient, together with naturality of the unit, gives
the stated inertia module.  Alternatively exactness extends this
natural comparison from the unit to its thick closure.  The two
branch coordinates may be inverse; their induced inertia map,
rather than a chosen identification of cohomology lines, records
that sign.

At each level retain the entire localization triangle and its
connecting arrow:
\[
 j_{n,!}V_n\longrightarrow Rj_{n,*}V_n
   \longrightarrow i_{n,*}L_{n,r}(V)
   \xrightarrow{\delta_n}j_{n,!}V_n[1].
\]
Open pullback and the adjunction comparison for $Rj_*$ give the
maps of these actual triangles; the connecting squares are obtained
from their cofiber diagrams.  Apply geometric Verdier duality to
those finite diagrams, retaining the smooth-branch dualizing factor
$\Lambda_r(1)[2]$ on the open side and point duality on the closed
side.  This gives the inverse diagrams of the actual adjunction
mates and their trace boundaries.  In particular on the closed
link they satisfy
\[
 R\varprojlim_n D_{\Lambda_r}L_{n,r}(V)
  \simeq D_{\Lambda_r}\operatorname*{hocolim}_nL_{n,r}(V),
\]
with the evaluation pairing.  For a constant coefficient the dual
of the degree-one transition is again multiplication by
$\ell^{m-n}$ and is pro-zero at fixed $r$; the surviving degree-zero
pairing is ordinary evaluation.  Finite source Verdier factors are
retained in the boundary diagram; this scalar link identity does
not remove them.  No equality $q^!=q^*$ for arbitrary coefficients
on a ramified normal pair is used.
\end{proof}

The lemma compares coefficient links and their dual boundaries for
the fixed branch charts.  It does not compute the product links of
the completed diagonal, normalized rectangles, or contraction graph
kernels.  Those product-link maps, involving distribution limits as
well as the continuous-action index, are required before a global
kernel functor can be glued.  In particular a coefficient branch
comparison cannot be used as a substitute for those kernel maps.

\subsection{One actual diagram for the four compatibility checks}
\label{sec:root-common-diagram}
We organize the remaining root comparisons by one diagram, with a
precise distinction between identities of assembled diagrams and
invertibility of their interpretation as global kernel operations.
This avoids assuming the desired coherence in a definition of the
root source.

At precision $r$, retain the actual finite neutralized root charts
of the four fixed models, their support products, their diagonal
and graph correspondences, and their finite Kummer branch charts.
For an expression in already specified finite coefficients, form
its relative bar and Kummer Cech diagrams on these actual products.
Include the open/closed localization arrows, the chosen carrier
projection, the two refinement paths, and the evaluation arrows.
Call this diagram $\mathcal B_r$.  Its maps are the actual pullback,
shriek direct image, units, counits and their Beck--Chevalley mates
in the fixed enhancement.  In particular the nonrepresentable
$!$ leg is not replaced by its $*$ leg.  No unspecified global
normalized coefficient is inserted into $\mathcal B_r$.

There is a distribution-presentation index $\mathbf n$ and a
continuous-action index $\mathbf m$ dominating it, one entry for
each inertia factor.  On discrete continuous mapping/cochain faces
the order is $\operatorname*{hocolim}_{\mathbf m}$ followed by
$R\varprojlim_{\mathbf n}$.  On their trace-dual faces use the dual
order prescribed by evaluation.  Derived precision recovery is
last.  Equal-variance indices may be synchronized by cofinality;
a direct and an inverse index may not be synchronized this way.

\begin{proposition}[Four faces of the same assembled diagram]
\label{prop:root-four-faces}
On the local completed kernel source and the coefficient branch
faces constructed above, the following comparisons come from
$\mathcal B_r$, and are compatible wherever their operations have
already been identified:
\begin{enumerate}
\item units and convolution substitution;
\item localization/cofiber connecting arrows;
\item evaluation, Verdier mates and carrier traces;
\item the actual refinement/base-change comparisons and their
composition simplices.
\end{enumerate}
Their common higher compatibility is inherited from the same bar,
adjunction and trace diagram.  For a global four-model root kernel
realization, this assertion becomes applicable only after the
product-link and completion comparison maps have been constructed
and proved invertible; it does not assert those missing maps.
\end{proposition}
\begin{proof}
For a string of actual correspondences, evaluate its coefficient
on the full product and use the multiple relative bar.  A
substitution deletes an intermediate bar by its augmentation;
inserting the diagonal gives its degeneracy and regular unit.
Whiskering two deletions in either order is the same two-dimensional
bar face.  This supplies the first comparison and its associativity
cells.  In the local completed module model these are exactly the
evaluations used in Theorem~\mainref{thm:completed-root-shriek}.

Retain an open/closed arrow in that same diagram instead of
collapsing it to its object values.  Its cofiber vertex and the
connecting arrow are then functorially determined.  Applying a
natural comparison $\tau:F\to G$ gives
$G(\delta)\tau_{Q_2}=\tau_{Q_0}[1]F(\delta)$, with its specified
homotopy, for the boundary $\delta:Q_2\to Q_0[1]$.
Finite cones and retracts, and every limit or colimit functor used
in the assembly, are exact.  Thus this identity survives assembly
on the specified diagrams.  Lemma~\ref{lem:root-branch-link}
identifies the actual single-support branch boundary face; no
identification of a product-kernel boundary follows just from it.

Taking a mate means applying the given adjunctions to this same
face.  Its evaluation satisfies, for an actual finite refinement,
\[
 e_n\circ(t_{m,n}\otimes1)
       =e_m\circ(1\otimes\operatorname{inf}_{n,m}).
\]
For composable structural maps the final pairing uses
$\operatorname{Tr}_{gf}
=\operatorname{Tr}_g\circ Rg_!\operatorname{Tr}_f$.
These are faces of the finite six-operation diagram, including
$!$ on nonproper projections.  The carrier comparison is precisely
this pairing followed by the source-augmentation evaluation of
Theorem~\mainref{thm:selected-carrier-interface}.  Consequently taking
mates, inserting a unit and contracting another kernel give the
same evaluation diagram.  Biduality is used on selected perfect
outputs or on the specified discrete/pseudocompact dual pair, not
on an unrestricted completed sheaf class.

Finally the refinement simplex is the composition simplex of the
actual chart maps and their adjunctions.  In local graph coordinates
it is $R_a\widehat\otimes_{R_b}R_b\to R_a$ and its iterated
substitution, as in Proposition~\mainref{prop:completed-root-graphs}.
Its base-change faces are the prescribed Beck--Chevalley maps,
not arbitrary identifications chosen after comparing stalks.
Taking a mate of this simplex or inserting the localization arrow
therefore gives the same simplex in the larger diagram.

Apply the assembly functors to $\mathcal B_r$ \emph{as a diagram}.
A functor carries simplices and their substitutions to simplices;
thus all four identities and their mixed cells remain identities.
This argument does not say that assembly commutes with an arbitrary
relative tensor or geometric direct image.  Those are additional
comparison maps.  They are proved for the local unit/rectangle
words and coefficient faces already used here.  For the still
missing product links they have not been proved.  Once an actual
comparison of diagrams has invertible components there as well,
restriction to the full generated source will retain the common
cells without four separate choices of coherence data.
\end{proof}


\begin{lemma}[Independent rectangular product-link faces]
\label{lem:root-rectangular-link-faces}
For two fixed branch charts and their punctured inclusions
$j_x,j_y$, use the product-open link
\[
 (i_x\times i_y)^*R(j_x\times j_y)_*(V\boxtimes W).
\]
This is the link to the product of the two punctured branches, not
the link to the full punctured neighbourhood of the corner.
Exterior products of actual finite, underlying-perfect selected
coefficients have this product-link comparison given by the tensor
product of the two branch bar diagrams.  Their two restriction orders agree, including the
localization-arrow faces.  Continuous-action colimits and the
dual trace diagrams retain this comparison.  With distribution
coefficients the same statement holds with independent presentation
indices and the completed exterior product, defined by their
inverse system; it does not exchange a distribution inverse limit
with a continuous-action colimit.
\end{lemma}
\begin{proof}
At fixed coefficient presentations the product Kummer torsor is
the product of the two actual torsors.  Its two Cech directions
form a bisimplicial diagram.  The scalar product-link map is the
external-product/Kunneth comparison on those nerves.  The branch
cohomology has finite amplitude and the coefficients here are
underlying-perfect, so totalization of this finite-dimensional
branch calculation gives the comparison with the product link.
The two orders are identified by the same bisimplicial Fubini map,
with the Koszul signs; this proves equality of the specified maps.
Tensor with, or restrict, the actual open/closed arrow before this
calculation to obtain the localization faces of the same diagram.

Take the continuous-action colimits at fixed presentations first.
Their tensor comparison is valid since derived tensor preserves
these colimits separately.  On the dual trace diagrams use derived
Hom of this equality, with their finite evaluation pairings.
For distribution coefficients retain the two presentation indices
and then take their inverse limit; this is the completed exterior
product, not an uncompleted tensor product of the two limits.
This proves precisely the claimed comparison.  It applies to
independent exterior factors and this product-open face.  Other
corner faces still have to be retained in a full gluing diagram.
A diagonal or graph imposes a
relation between the two inertia directions and is not such a
factorization.  Nor does this lemma construct the missing global
normalized source factor on a whole refined fibre.
\end{proof}


\section{Finite unit and dual-unit diagrams}
\begin{lemma}[Unit and dual-unit diagrams on the complete supports]
\label{lem:root-unit-dual-unit-diagram}
At finite root precision and level, the preceding adjunction
construction applies to every actual transition $q:Y\to X$
between the fixed support charts, including charts meeting more
than one branch.  With $F_X=\Delta_{X,!}\Lambda_r$ it gives
\[
 (q\times q)^*F_X\longrightarrow F_Y,
 \qquad
 (q\times q)_!\Delta_{Y,*}\omega_Y
                  \longrightarrow\Delta_{X,*}\omega_X,
 \quad \omega_X=D_X\Lambda_r.
\]
Composition, restriction of the localization arrow, and the
Verdier-mate identities belong to the same actual adjunction
diagram.  These maps are not asserted invertible.  In particular,
a refinement does not identify its fibre relation with its diagonal.
The second family is a family of \emph{dual-unit} trace maps; on a
singular support it is not a family of shriek-unit maps obtained
by cancelling one constant twist and shift.
\end{lemma}
\begin{proof}
The construction of $\beta$ in
Proposition~\mainref{prop:root-diagonal-branch-transition} only uses the
actual diagonal counit and composition of pullbacks.  It does not
use smoothness.  Diagonals of the finite separated root charts here
are proper and representable.  Proper duality therefore identifies
$D F_X=\Delta_{X,*}\omega_X$, giving the second family by taking
the mate of the first.  The unit/counit identities for three
composable maps identify the two resulting $\beta$ maps.  Their
dual identity is trace transitivity.  Retaining the open/closed
arrow and its cofiber in this natural transformation gives its
localization and boundary faces.  This is a particular actual
instance of the diagram in Proposition~\ref{prop:root-four-faces};
it requires no choice of a new comparison on each branch.  For a
refinement, the first source still has support on the full inverse
image of the diagonal.  Nothing in the counit construction asserts
that this support equals the fine diagonal or that the counit is
an equivalence.  Finally proper duality gives $\omega_X$ itself,
not an invertible replacement for it.
\end{proof}

\section*{Acknowledgments}

This mathematical paper was prepared with the assistance of artificial intelligence.  I
first acknowledge the textual data produced by humanity as a whole, which
provides the linguistic and cultural record from which mathematical expression
can be learned.  I then thank the accumulated scholarly achievements of
mathematicians throughout history.  I also thank the researchers whose work has
advanced and refined neural networks through successive generations.  Finally,
I thank the OpenAI team for the training and development of the models used in
the preparation of this paper.

\begin{thebibliography}{9}
\bibitem{Local} Y.~Jiang,
\emph{Local Kummer Realizations and Verdier Duality},
local submission source, 22 September 2026.
\bibitem{Refinement}
\emph{Morita Coefficients, Incidence, and Semistable Refinement},
local source consulted 22 September 2026.
\bibitem{LZEnhanced} Y.~Liu and W.~Zheng,
\emph{Enhanced six operations and base change theorem for higher
Artin stacks}, arXiv:1211.5948.
\bibitem{LZAdic} Y.~Liu and W.~Zheng,
\emph{Enhanced adic formalism and perverse t-structures for higher
Artin stacks}, arXiv:1404.1128.
\bibitem{Nearby}
\emph{Verdier-Compatible Kernel Skeletonization under Anchored Semistable
Refinement}, nearby theorem and local-root appendices, 27 September 2026;
source in \texttt{../verdier\_kernel\_skeletonization/}.
\end{thebibliography}
\end{document}
